% =============================================================
%  Module Description — Software Engineering (CS303-F28)
%  Shown on https://wosm.academy as the course page for CS303-F28.
%  Compile with: pdflatex cs303-module-description.tex
%  Write "TODO(owner): ..." where information is still missing;
%  the website highlights it.
% =============================================================
\documentclass[11pt,a4paper]{article}

\usepackage[utf8]{inputenc}
\usepackage[T1]{fontenc}
\usepackage{lmodern}
\usepackage[margin=2.2cm]{geometry}
\usepackage{booktabs}
\usepackage{tabularx}
\usepackage{array}
\usepackage{enumitem}
\usepackage{xcolor}
\usepackage{titlesec}
\usepackage{fancyhdr}
\usepackage[hidelinks]{hyperref}

% ---------- Module details (edit these) ----------------------
% \ModuleCode must equal the course's `code` in src/content/courses/.
\newcommand{\ModuleTitle}{Software Engineering}
\newcommand{\ModuleCode}{CS303-F28}
\newcommand{\Programme}{Foundation programme}
\newcommand{\Level}{Semester 5}
\newcommand{\Semester}{Fall 2028}
\newcommand{\Credits}{8 ECTS}
\newcommand{\ModuleLeader}{Dr. Ali Khudiyev}
\newcommand{\Email}{ali.khudiyev@wosm.academy}
\newcommand{\OfficeHours}{Monday--Friday, 10:00--17:00 (Baku time, UTC+4)}
\newcommand{\Prerequisites}{CS101-F26 Introduction to Informatics; CS103-S27 Data Structures and Algorithms I}

% ---------- Styling ------------------------------------------
\definecolor{accent}{RGB}{0,70,130}
\titleformat{\section}{\large\bfseries\color{accent}}{\thesection.}{0.5em}{}
\titlespacing*{\section}{0pt}{1.2em}{0.5em}
\setlist{itemsep=2pt, topsep=4pt}
\setlength{\parindent}{0pt}
\setlength{\parskip}{4pt}

\pagestyle{fancy}
\fancyhf{}
\lhead{\small \ModuleCode\ -- \ModuleTitle}
\rhead{\small Module Description}
\cfoot{\small \thepage}
\renewcommand{\headrulewidth}{0.4pt}

% =============================================================
\begin{document}

\begin{center}
  {\LARGE\bfseries\color{accent} \ModuleTitle}\\[4pt]
  {\large Module Description}
\end{center}
\vspace{6pt}

% ---------- Summary table (not shown on the website: the course page
% shows these details from the macros above) --------------------------
\begin{tabularx}{\textwidth}{>{\bfseries}l X >{\bfseries}l X}
\toprule
Module code   & \ModuleCode    & Credits      & \Credits \\
Programme     & \Programme     & Level        & \Level \\
Semester      & \Semester      & Prerequisites & \Prerequisites \\
Module leader & \ModuleLeader  & Email        & \href{mailto:\Email}{\Email} \\
Office hours  & \multicolumn{3}{l}{\OfficeHours} \\
\bottomrule
\end{tabularx}

% ---------- Sections (shown on the website) ------------------
\section{Module Overview}
Writing a program that works is different from building software that a team can
develop, test, ship and maintain for years. Software Engineering is about the second
problem. It covers how software projects are organised, how requirements are
captured, how systems are designed and kept maintainable, and how quality is
assured through testing and review.

The module is practical. Alongside the lectures you will work in a small team on a
software project in Java, using the tools and practices of professional development:
version control, automated tests, continuous integration and code review. The module
closes the core of the Foundation programme and prepares you for the optional
research or project work that follows Semester 5.

\section{Aims}
\begin{itemize}
  \item Introduce the processes, methods and tools used to develop software systematically.
  \item Develop your ability to design software that is easy to understand, test and change.
  \item Teach systematic testing and other quality-assurance practices.
  \item Give experience of working in a team on a substantial software project.
\end{itemize}

\section{Learning Outcomes}
After completing this course, you will be able to:
\begin{enumerate}[label=LO\arabic*.]
  \item Compare software process models, including plan-driven and agile approaches such as Scrum, and choose one suited to a project.
  \item Elicit and specify requirements as user stories, use cases and acceptance criteria.
  \item Design software with UML models, design principles such as low coupling, high cohesion and SOLID, and design patterns, and improve existing code by refactoring.
  \item Test software systematically with unit, integration and system tests, use test-driven development, and measure test coverage.
  \item Use version control, build automation, continuous integration and code review in a team workflow.
  \item Plan and deliver a software project in a team, and discuss the professional, ethical and legal responsibilities of software engineers, including open-source licensing.
\end{enumerate}

\section{Syllabus}
\begin{tabularx}{\textwidth}{c X}
\toprule
\textbf{Unit} & \textbf{Topic} \\
\midrule
1  & What is software engineering? Software processes and life cycles \\
2  & Agile methods: Scrum and Kanban \\
3  & Requirements engineering: elicitation, user stories and use cases \\
4  & Version control with Git and collaborative workflows \\
5  & Software modelling with UML \\
6  & Design principles: abstraction, coupling and cohesion, SOLID \\
7  & Design patterns I: creational and structural patterns \\
8  & Design patterns II: behavioural patterns; software architecture styles \\
9  & Clean code and refactoring \\
10 & Software testing I: unit testing and test-driven development \\
11 & Software testing II: integration, system and acceptance testing; coverage \\
12 & Build automation, continuous integration and continuous delivery \\
13 & Code review, static analysis and software quality \\
14 & Software maintenance and evolution; technical debt \\
15 & Project management: estimation, planning and risk \\
16 & Professional practice, ethics and open-source licensing \\
17 & Project presentations \& Course Review \\
\bottomrule
\end{tabularx}

\section{Teaching and Learning Methods}
The module is self-paced. All materials are provided on Moodle.
Each unit has a lecture and exercises. Throughout the module you work in a small team
on a software project, applying the practices of each unit as you go.

\begin{tabularx}{0.6\textwidth}{X r}
\toprule
\textbf{Activity} & \textbf{Hours} \\
\midrule
Lectures            & 30 \\
Labs / Tutorials    & 50 \\
Independent study   & 150 \\
\midrule
\textbf{Total}      & \textbf{230} \\
\bottomrule
\end{tabularx}

\section{Assessment}
Your knowledge will be assessed with quizzes, the team project and the final exam.
The project is assessed on the delivered software and on its requirements, design,
tests and development history, together with a short individual reflection. The exam
checks whether you can apply design principles, patterns and testing techniques to
given problems.

\begin{tabularx}{\textwidth}{X c c c}
\toprule
\textbf{Component} & \textbf{Weight} & \textbf{Due} & \textbf{LOs} \\
\midrule
Quizzes & 15\% & Throughout & LO1--LO4 \\
Team project and individual reflection & 50\% & Throughout & LO2--LO6 \\
Final exam & 35\% & End of module & LO1--LO4, LO6 \\
\midrule
\textbf{Total} & \textbf{100\%} & & \\
\bottomrule
\end{tabularx}

\section{Reading List}
\textbf{Essential}
\begin{itemize}
  \item I.\ Sommerville: \textit{Software Engineering}, 10th ed., Pearson, 2016.
\end{itemize}
\textbf{Recommended}
\begin{itemize}
  \item E.\ Gamma, R.\ Helm, R.\ Johnson, J.\ Vlissides: \textit{Design Patterns: Elements of Reusable Object-Oriented Software}, Addison-Wesley, 1994.
  \item M.\ Fowler: \textit{Refactoring: Improving the Design of Existing Code}, 2nd ed., Addison-Wesley, 2018.
  \item J.\ Bloch: \textit{Effective Java}, 3rd ed., Addison-Wesley, 2018.
  \item D.\ Thomas, A.\ Hunt: \textit{The Pragmatic Programmer}, 20th anniversary ed., Addison-Wesley, 2019.
  \item T.\ Winters, T.\ Manshreck, H.\ Wright: \textit{Software Engineering at Google}, O'Reilly, 2020.
\end{itemize}

\section{Materials and Licence}
Course materials are published under the
\href{https://creativecommons.org/licenses/by-sa/4.0/}{Creative Commons Attribution-ShareAlike 4.0}
licence unless stated otherwise. See \href{https://wosm.academy/licence/}{wosm.academy/licence} for details.

\section{Policies}
\begin{itemize}
  \item \textbf{Attendance:} Attendance is not mandatory. The module is self-paced, so you can study the materials when it suits you.
  \item \textbf{Late submissions:} Work submitted after the deadline is penalised.
  \item \textbf{Academic integrity:} All submitted work must be your own.
        See the \href{https://wosm.academy/code-of-conduct/}{Code of Conduct}.
\end{itemize}

\end{document}
